\chapter{Modes of Access as Task-Relative, Not Object-Relative}
\label{ch:mode-relativity}

\begin{quotation}
\noindent\textit{Synopsis.} This chapter synthesizes the preceding
discussion by arguing that access mode is determined by the relationship
among consumer, task, and time, not by an intrinsic property of an
artifact. The same object can function as a preserved source for one task
while acting as an incomplete trace for another; proofs, paintings,
contracts, and memories are revisited to show how a single artifact
occupies different points on the fidelity continuum depending on what
continuation is being attempted. This dissolves the earlier
preservation/reconstitution binary into a more general relational
framework.
\end{quotation}

\section{The Criterion, Made Precise}

Chapter~\ref{ch:fidelity-continuum} defined fidelity for a rendering $r$
of an object $e$ relative to a task $\tau$. What was left implicit there
is that $r$ itself depends on who is asking and when: the substrate a
particular consumer actually has in hand is not a fixed fact about the
world, it is a fact about that consumer at that moment.

\begin{definition}[Substrate Available to a Consumer]
\label{def:available-substrate}
For a consumer $c$ at time $t$, $r_{c,t}$ denotes the rendering
(\cref{def:rendering}) $c$ has causal access to at $t$.
\end{definition}

\begin{definition}[Access Mode]
\label{def:access-mode}
The access mode of consumer $c$ to a task $\tau$ at time $t$ is
$M(c,\tau,t) := \fidelity{r_{c,t}}{\tau}$.
\end{definition}

Under this formalization, \cref{def:access-mode} is not an independent
postulate about the nature of access; it is fidelity, already defined,
simply evaluated at whichever substrate a specific consumer's specific
circumstances make available. This is a stronger and more useful move
than describing mode as ``a relation among consumer, task, and time'' in
the abstract, because it says exactly what that relation \emph{is}: not
a new primitive, but an old one indexed correctly.

\begin{proposition}[Mode as a Relation]
\label{prop:mode-relation}
Access mode is a relation $M(c, \tau, t)$ over consumer $c$, task $\tau$,
and time $t$ -- not a property $M(A)$ of an artifact $A$ in isolation.
The same artifact can occupy different points on the fidelity continuum
of \cref{def:fidelity} simultaneously for different consumer-task pairs.
\end{proposition}

\begin{proofsketch}
Immediate from \cref{def:access-mode}: $M(c,\tau,t)$ is defined as a
function of $r_{c,t}$ and $\tau$ jointly, and $r_{c,t}$ is itself
consumer- and time-indexed by \cref{def:available-substrate}. $M$
collapses to a property of a single artifact $A$ alone only in the
special, non-generic case where $r_{c,t} = A$ for every consumer and
every time -- which is exactly what fails whenever different consumers
have access to different renderings of the same underlying object, or
the same consumer's access changes over time.
\end{proofsketch}

\section{Same Artifact, Different Modes}

A formal proof is high-fidelity (\cref{def:fidelity}) for the task of
verification, since $r_{c,t}$ for a verifier with the text in hand is
close to the proof itself; the same proof is lower-fidelity for the task
of explanation or generalization, since $r_{c,t}$ for a reader attempting
those tasks is typically their memory or partial recollection of the
argument, not the text. A painting is, for its painter during production,
close to direct access; for a later viewer inferring causal history from
static cues (Chapter~\ref{ch:painting}), two separate rendering steps are
at issue and are worth keeping apart. The step from event to canvas is
already low-fidelity for the causal-inference task, by
\cref{prop:open-causal-space} alone, regardless of who views the canvas
or how directly they perceive it. The step from canvas to a specific
viewer's percept is the one Chapter~\ref{ch:mem8}'s fork actually bears
on (\cref{oq:mem8-fork}): how close $r_{c,t}$ sits to the canvas itself.
By \cref{prop:admissibility-monotone}, the viewer's fidelity for the
causal-inference task is bounded above by the canvas's own already-low
fidelity for that task, so resolving the fork can only leave the
viewer's fidelity the same or lower -- it cannot resolve the ambiguity
\cref{prop:open-causal-space} already established into certainty. The
benign-ambiguity conclusion of \cref{prop:benign-ambiguity} is therefore
fork-independent, even though the raw perceptual fidelity of viewing the
canvas is not: the proof of \cref{prop:benign-ambiguity} requires only
that the task tolerate multiple continuations, not that the viewer's
perception of the canvas be unmediated. A contract clause is
high-fidelity for the task of reading its literal text; the same
contract is low-fidelity for the task of establishing original intent,
which is typically litigated by reconstructing a trace-based
approximation of a past state of mind (Chapter~\ref{ch:contracts}) --
$r_{c,t}$ for a court asking about intent is correspondence, drafting
history, and testimony, not the contract text itself.

\begin{center}
\small
\begin{tabular}{p{1.4cm}p{2.0cm}p{2.6cm}p{2.6cm}}
\toprule
Artifact & Task & $r_{c,t}$ & Fidelity \\
\midrule
Proof & Verification & The proof text itself & High \\
Proof & Explanation & Reader's recollection of the argument & Low \\
Painting & Production & Canvas under the painter's hand & High \\
Painting & Causal inference & The canvas itself & Low, fork-independent \\
Painting & Perceiving the canvas & Viewer's percept of the canvas & Open (\cref{oq:mem8-fork}) \\
Contract & Literal reading & The contract text itself & High \\
Contract & Establishing intent & Correspondence, drafting history, testimony & Low \\
\bottomrule
\end{tabular}
\end{center}

The table's third column is doing the real work: in every row, what
changes across tasks for the same artifact is not the artifact but which
substrate a consumer pursuing that particular task actually has access
to, exactly as \cref{def:access-mode} formalizes.

\section{Why This Dissolves Rather Than Answers the Earlier Question}

It is worth being precise about what this chapter, and
Part~\ref{part:modes-of-access} as a whole, has and has not accomplished.
It has not settled whether embodied perception ever achieves fidelity
exactly $1$; that question is left explicitly to
Chapter~\ref{ch:mem8}'s unresolved fork. What it has done is show that
the original question motivating this part of the monograph --- is a
given artifact accessed in preservation mode or reconstitution mode? ---
was not well posed to begin with, because it asked for a property of the
artifact when the answer was always a property of the artifact together
with a consumer, a task, and a time. Once \cref{def:access-mode} makes
that dependency explicit, the preservation/reconstitution distinction
does not need to be defended as a real boundary somewhere on the fidelity
continuum; it can simply be retired in favor of the continuum itself,
with ``preservation'' surviving only as a convenient name for the
high-fidelity end of a scale that has no other privileged points.
Part~\ref{part:worked-domains} proceeds on this basis: rather than asking
of chess, painting, maps, proofs, contracts, science, memory, and media
pipelines which mode of access each one exemplifies, it asks, for each
domain and each task, what $r_{c,t}$ actually is and how close it sits to
the source.
